%% ═══════════════════════════════════════════════════════════════════════════
%% CHAPTER 6: COMPOSITING, POST-PRODUCTION & VIDEO EDITING
%% Thicc Splines Save Lives — A Comprehensive Blender User Manual
%% ═══════════════════════════════════════════════════════════════════════════

\chapter{Compositing, Post-Production \& Video Editing}
\label{ch:compositing}

\begin{quotebox}
\itshape ``The art of compositing is the art of the invisible join. When it works, the audience never knows where the real ends and the synthetic begins.''

\upshape --- John Knoll, co-creator of Adobe Photoshop and Visual Effects Supervisor, ILM
\end{quotebox}

\vspace{1em}
Rendering is not the final step. In any professional pipeline --- from a 30-second social media clip to a feature film --- the raw render is merely the starting point for a process of refinement, correction, and integration that can transform a competent image into a compelling one. This chapter covers the full post-production pipeline available within Blender: the node-based Compositor for image processing and color grading, the Video Sequence Editor for non-linear editing, motion tracking for VFX integration, and Grease Pencil for 2D/3D hybrid work. Every tool discussed here operates within a single Blender file, meaning you can move from render to final cut without leaving the application.

\begin{historynote}
The concept of compositing predates digital tools by decades. The \textbf{optical printer}, invented in the 1920s and refined through the 1970s, was a mechanical device that re-photographed film through multiple exposures, combining foreground and background elements frame by frame. Linwood Dunn at RKO perfected the technique for films like \textit{King Kong} (1933) and \textit{Citizen Kane} (1941). The process was painstaking: each composite might require dozens of film passes, each introducing grain, color shift, and registration errors.

The digital revolution arrived in the late 1980s when Industrial Light \& Magic (ILM) developed early digital compositing for \textit{The Abyss} (1989) and \textit{Terminator 2: Judgment Day} (1991). Apple's Shake (acquired from Nothing Real in 2002) and The Foundry's Nuke became the industry-standard node-based compositors. Blender's Compositor is a direct descendant of this lineage --- a node-based, non-destructive image processing system designed for the same workflows that professionals use in Nuke and Fusion, scaled to run on a single workstation.
\end{historynote}

%% ═══════════════════════════════════════════════════════════════════════════
\section{Compositor Workspace and Architecture}
\label{sec:comp-architecture}
%% ═══════════════════════════════════════════════════════════════════════════

\subsection{Overview}

Blender's \textbf{Compositor}\index{Compositor} is a node-based image processing system that operates on rendered images after the 3D render is complete. It allows post-production operations --- color correction, effects, compositing of multiple render layers, mixing CG with live footage --- entirely within Blender, without exporting to external software. For artists coming from a Photoshop workflow, think of the Compositor as a non-destructive, procedural version of adjustment layers and blend modes, but with the full power of a node graph.

\subsection{Accessing the Compositor}

There are two ways to access the Compositor:

\begin{itemize}[leftmargin=2em]
  \item Switch to the \textbf{Compositing} workspace via the header tabs, which provides a pre-configured layout with the Compositor editor and an Image Editor for preview.
  \item Change any editor area to the Compositor editor type manually.
\end{itemize}

Two checkboxes in the Compositor header are critical:

\begin{itemize}[leftmargin=2em]
  \item \textbf{Use Nodes} --- activates node-based compositing. Without this checkbox, the raw render passes directly to the output with no processing.
  \item \textbf{Backdrop} --- displays the compositing result directly in the node editor background, providing immediate visual feedback as you adjust nodes.
\end{itemize}

\subsection{Node Tree Architecture}

The Compositor uses a \textbf{node tree}\index{node tree} organized as a directed acyclic graph (DAG) where data flows left-to-right through connections called \textit{noodles}. Nodes fall into three categories:

\begin{itemize}[leftmargin=2em]
  \item \textbf{Input nodes} provide data: Render Layers, Image, Movie Clip, Value, RGB.
  \item \textbf{Processing nodes} transform data: filters, color adjustments, masks, distortions.
  \item \textbf{Output nodes} display or save results: Composite, Viewer, File Output, Split Viewer.
\end{itemize}

Each connection carries one of several data types, indicated by wire color:

\vspace{0.5em}
\rowcolors{2}{rowgray}{white}
\begin{tabularx}{\textwidth}{l l X}
\toprule
\rowcolor{primary}
\tableheader{Data Type} & \tableheader{Wire Color} & \tableheader{Description} \\
\midrule
Color (RGBA) & Yellow & Standard 4-channel color image \\
Value (float) & Grey & Single-channel grayscale data (depth, alpha, factors) \\
Vector & Blue & Multi-component data (normals, UV coordinates, speed vectors) \\
\bottomrule
\end{tabularx}
\vspace{0.5em}

\begin{tipbox}
When connecting nodes, Blender will automatically convert between data types where possible --- a Color output connected to a Value input will be converted to grayscale. However, connecting a Value output to a Vector input may produce unexpected results. Pay attention to wire colors to ensure your data flows are semantically correct.
\end{tipbox}

\subsection{Real-Time Compositor}
\label{subsec:realtime-compositor}

Blender 3.5 introduced a \textbf{Real-Time Compositor}\index{Real-Time Compositor} that evaluates node trees in the viewport with GPU acceleration. As of Blender 4.4, the real-time compositor supports most common nodes, though some complex nodes (like Glare with certain modes) may fall back to CPU processing.

To enable the real-time compositor: \textbf{Render Properties} $\rightarrow$ enable the \textbf{Compositor} checkbox $\rightarrow$ set Mode to \textbf{Realtime}.

This feature is transformative for iterative color grading --- you can adjust a Color Balance node and see the result update in the viewport in real time, rather than waiting for a full re-render. For artists accustomed to working in Photoshop or DaVinci Resolve where adjustments are immediate, the real-time compositor brings Blender's compositing workflow much closer to that interactive experience.

\subsection{Essential Keyboard Shortcuts}

\begin{shortcutbox}
\textbf{Compositor Keyboard Shortcuts}

\vspace{0.3em}
\rowcolors{2}{rowgray}{white}
\begin{tabularx}{\textwidth}{l X}
\toprule
\rowcolor{primary}
\tableheader{Shortcut} & \tableheader{Action} \\
\midrule
Shift+A & Add Node menu \\
Ctrl+Shift+Click & Connect node to Viewer for preview \\
M & Mute selected node (bypasses it in the chain) \\
H & Hide unused sockets (collapses the node visually) \\
Ctrl+J & Join selected nodes into a Frame (organizational) \\
Ctrl+G & Group selected nodes into a reusable Node Group \\
Tab & Enter/exit a Node Group for internal editing \\
Ctrl+H & Hide/show all unconnected sockets \\
Backspace & Recalculate (when a Viewer is connected) \\
\bottomrule
\end{tabularx}
\end{shortcutbox}

%% ═══════════════════════════════════════════════════════════════════════════
\section{Render Passes}
\label{sec:render-passes}
%% ═══════════════════════════════════════════════════════════════════════════

\subsection{What Are Render Passes?}

A \textbf{Render Pass}\index{render pass} isolates a specific component of the rendered image --- diffuse color, specular highlights, shadows, depth, and so on. By rendering these components separately, compositors can adjust each independently and recombine them for maximum control. This is the foundation of professional compositing: rather than trying to fix a finished image, you work with its constituent parts.

Render passes are enabled in \textbf{View Layer Properties} $\rightarrow$ \textbf{Passes} and appear as individual outputs on the \textbf{Render Layers} node in the Compositor. The more passes you enable, the more control you have in post-production --- but each pass adds to render time and file size, so enable only what you need.

\subsection{Light Passes}

Light passes separate the contribution of different shading components. In Cycles, the combined image is the mathematical sum of all light passes, which means you can decompose and recompose the image perfectly:

\vspace{0.5em}
\rowcolors{2}{rowgray}{white}
\begin{tabularx}{\textwidth}{l L{4cm} X}
\toprule
\rowcolor{primary}
\tableheader{Pass} & \tableheader{Description} & \tableheader{Typical Use} \\
\midrule
Diffuse Light & Intensity and color of light hitting Diffuse/SSS BSDFs (excludes surface color) & Brighten/darken shadows without affecting material colors \\
Diffuse Color & Surface colors of Diffuse/SSS BSDFs (excludes lighting) & Adjust material colors in compositing \\
Glossy Light & Light reflected by Glossy BSDFs & Control specular highlight intensity \\
Glossy Color & Surface color of Glossy BSDFs & Tint or desaturate reflections \\
Transmission Light & Light passing through Transmission BSDFs (glass, translucency) & Control glass brightness \\
Transmission Color & Surface color of Transmission BSDFs & Tint glass colors \\
Volume Light & Light contribution from volume shaders & Control fog/volumetric brightness \\
Emission & Direct emission from surfaces (emissive materials, lights via shader) & Add bloom selectively to emissive objects \\
\bottomrule
\end{tabularx}
\vspace{0.5em}

The recombination formula is exact:

\begin{codebox}
Combined = (Diffuse Light * Diffuse Color)\\
\hspace*{2em}     + (Glossy Light * Glossy Color)\\
\hspace*{2em}     + (Transmission Light * Transmission Color)\\
\hspace*{2em}     + Volume Light\\
\hspace*{2em}     + Emission
\end{codebox}

\begin{tipbox}
Understanding this formula is key to effective compositing. If you want to make reflections brighter without affecting diffuse lighting, multiply the Glossy Light pass by a value greater than 1.0 before recombining. If you want to change the color of glass without re-rendering, adjust the Transmission Color pass. The formula guarantees that your adjusted result will be physically consistent with the original lighting.
\end{tipbox}

\subsection{Data Passes}

Data passes carry non-color information about the scene geometry:

\vspace{0.5em}
\rowcolors{2}{rowgray}{white}
\begin{tabularx}{\textwidth}{l l L{3cm} X}
\toprule
\rowcolor{primary}
\tableheader{Pass} & \tableheader{Data Type} & \tableheader{Description} & \tableheader{Typical Use} \\
\midrule
Depth (Z) & Value (float) & Distance from camera to surface per pixel & Depth of field in compositing, fog, distance-based effects \\
Normal & Vector (XYZ) & World-space surface normals per pixel & Relighting, normal-based masking, edge detection \\
Mist & Value (0--1) & Linear depth ramp based on Mist Pass settings in World Properties & Distance fog, aerial perspective \\
Shadow & Value & Shadows collected by shadow catcher objects & Compositing CG shadows onto live footage \\
Ambient Occlusion & Value (0--1) & Proximity-based self-shadowing per pixel & Enhance or reduce crevice darkening \\
Denoising Data & Special & Albedo and Normal data used by the denoiser & Required for the Denoise node to function \\
\bottomrule
\end{tabularx}
\vspace{0.5em}

\subsection{Cryptomatte Passes}
\label{subsec:cryptomatte-passes}

\textbf{Cryptomatte}\index{Cryptomatte} is a modern matte extraction system developed by Psyop and adopted across the industry. It stores per-pixel ID information using a hashing scheme that preserves anti-aliased edges, motion blur, and transparency --- none of which are possible with the older Object Index system.

\vspace{0.5em}
\rowcolors{2}{rowgray}{white}
\begin{tabularx}{\textwidth}{l X}
\toprule
\rowcolor{primary}
\tableheader{Cryptomatte Pass} & \tableheader{Description} \\
\midrule
Object & Per-object ID matte with anti-aliased edges, motion blur, and transparency support \\
Material & Per-material ID matte \\
Asset & Per-asset ID matte \\
\bottomrule
\end{tabularx}
\vspace{0.5em}

Cryptomatte passes are used with the \textbf{Cryptomatte node} in the Compositor to create perfect mattes for individual objects or materials without requiring any pre-scene setup --- no manual Object Index assignments needed. See Section~\ref{subsec:cryptomatte-node} for the workflow.

\subsection{Index Passes (Legacy)}

The older masking system uses integer IDs assigned manually to objects and materials:

\vspace{0.5em}
\rowcolors{2}{rowgray}{white}
\begin{tabularx}{\textwidth}{l X L{4cm}}
\toprule
\rowcolor{primary}
\tableheader{Pass} & \tableheader{Description} & \tableheader{Setup Required} \\
\midrule
Object Index & Integer ID per object & Must set Pass Index on each object in Object Properties \\
Material Index & Integer ID per material & Must set Pass Index on each material in Material Properties \\
\bottomrule
\end{tabularx}
\vspace{0.5em}

These passes produce hard edges (no anti-aliasing) and do not support motion blur or transparency in the mask. They are largely superseded by Cryptomatte but remain useful for simple ID-based operations where setup speed matters more than edge quality.

\subsection{Passes in Eevee vs.\ Cycles}

Both render engines support render passes, but with important differences:

\vspace{0.5em}
\rowcolors{2}{rowgray}{white}
\begin{tabularx}{\textwidth}{l L{4cm} L{4cm}}
\toprule
\rowcolor{primary}
\tableheader{Feature} & \tableheader{Cycles} & \tableheader{Eevee} \\
\midrule
Diffuse/Glossy/Transmission split & Full separation & Available (simplified model) \\
Cryptomatte & Object, Material, Asset & Object, Material, Asset \\
Shadow Catcher & Yes & Yes \\
Depth pass & Raytraced (accurate) & Rasterized (Z-buffer) \\
Denoising data & Yes (for NLM/OptiX/OIDN) & Not needed (no path tracing noise) \\
AO pass & Raytraced & Screen-space \\
\bottomrule
\end{tabularx}
\vspace{0.5em}

\begin{warningbox}
Eevee's depth pass comes from the rasterized Z-buffer, which can have precision issues at large distances. If you plan to use depth-based compositing effects (fog, DOF) with Eevee renders, test the depth range carefully and consider using the Mist pass instead, which provides a normalized 0--1 range regardless of scene scale.
\end{warningbox}

%% ═══════════════════════════════════════════════════════════════════════════
\section{Multi-Layer EXR Workflow}
\label{sec:multilayer-exr}
%% ═══════════════════════════════════════════════════════════════════════════

\subsection{What Is Multi-Layer EXR?}

\textbf{OpenEXR Multi-Layer}\index{OpenEXR}\index{Multi-Layer EXR} is a high dynamic range image format that stores multiple render passes in a single file. Each pass becomes a named layer within the EXR container, preserving full 32-bit floating-point precision per channel. Developed originally at ILM for internal production use, OpenEXR is now the universal interchange format for visual effects work.

\subsection{Why Use Multi-Layer EXR?}

\vspace{0.5em}
\rowcolors{2}{rowgray}{white}
\begin{tabularx}{\textwidth}{l X}
\toprule
\rowcolor{primary}
\tableheader{Benefit} & \tableheader{Description} \\
\midrule
Non-destructive & All passes preserved; you can change compositing decisions without re-rendering \\
Full dynamic range & 32-bit float data preserves highlights and shadows that 8-bit formats clip \\
Single file & All passes in one file instead of dozens of separate images \\
Industry standard & Compatible with Nuke, After Effects (via plugin), Fusion, Natron \\
Re-renderable & If one pass looks wrong, you can isolate which component needs fixing \\
\bottomrule
\end{tabularx}
\vspace{0.5em}

\begin{tipbox}
The single most impactful workflow improvement for intermediate Blender users is adopting Multi-Layer EXR as the default render output format. It costs nothing in render time (the passes are computed either way), eliminates the need to re-render when you want to tweak compositing, and preserves the full dynamic range that 8-bit PNG or JPEG would destroy.
\end{tipbox}

\subsection{Exporting Multi-Layer EXR from Blender}

\begin{enumerate}[leftmargin=2em]
  \item \textbf{Output Properties} $\rightarrow$ \textbf{Output}
  \begin{itemize}[leftmargin=2em]
    \item Set File Format to \textbf{OpenEXR MultiLayer}
    \item Set Color Depth to \textbf{Float (Full)} for maximum precision
    \item Set Codec to \textbf{DWAA} (lossy but small), \textbf{ZIP} (lossless), or \textbf{PIZ} (lossless, good for noisy data)
  \end{itemize}
  \item \textbf{View Layer Properties} $\rightarrow$ \textbf{Passes} --- enable all desired passes (Diffuse Color, Diffuse Light, Glossy Light, Depth, etc.)
  \item \textbf{Render} (F12 for single frame, Ctrl+F12 for animation)
  \item Files are saved to the output path with all enabled passes embedded in a single EXR per frame
\end{enumerate}

\subsection{Importing Multi-Layer EXR into the Compositor}

\begin{enumerate}[leftmargin=2em]
  \item Add an \textbf{Image node} (Shift+A $\rightarrow$ Input $\rightarrow$ Image)
  \item Open the multi-layer EXR file
  \item The Image node displays all stored passes as individual outputs --- each output socket corresponds to one named layer in the EXR
  \item Connect passes to processing nodes as needed
\end{enumerate}

\subsection{File Output Node for Multi-Pass Export}

The \textbf{File Output}\index{File Output node} node provides more control than the standard output path:

\begin{enumerate}[leftmargin=2em]
  \item Add a \textbf{File Output} node (Shift+A $\rightarrow$ Output $\rightarrow$ File Output)
  \item In the node's properties sidebar (N-panel), set format to \textbf{OpenEXR MultiLayer}
  \item Add input slots for each pass you want to save
  \item Connect render passes to the corresponding inputs
  \item Each input becomes a named layer in the output EXR
\end{enumerate}

This is particularly useful for saving \textit{compositor-processed} passes --- for example, denoised passes or recolored passes --- back to EXR for use in external software like Nuke or DaVinci Resolve.

\begin{historynote}
The OpenEXR format was developed at Industrial Light \& Magic beginning in 1999 and open-sourced in 2003. The name ``EXR'' derives from the file extension, which itself was chosen because ILM's internal format was called ``EXR'' (for ``extended range''). Before EXR, VFX facilities used proprietary HDR formats that were incompatible across studios. EXR's open-source release unified the industry around a single high-dynamic-range format --- one of the most successful standardization efforts in visual effects history.
\end{historynote}

%% ═══════════════════════════════════════════════════════════════════════════
\section{Key Compositor Nodes}
\label{sec:compositor-nodes}
%% ═══════════════════════════════════════════════════════════════════════════

The Compositor provides dozens of nodes organized into categories. This section covers the nodes you will use most frequently, with enough detail to apply them effectively.

\subsection{Color Adjustment Nodes}

\subsubsection{Color Balance (Lift/Gamma/Gain)}
\label{subsec:color-balance}

The \textbf{Color Balance}\index{Color Balance node} node provides three-way color correction --- the industry-standard method for color grading used in every professional compositing and color grading application:

\vspace{0.5em}
\rowcolors{2}{rowgray}{white}
\begin{tabularx}{\textwidth}{l l X}
\toprule
\rowcolor{primary}
\tableheader{Control} & \tableheader{Affects} & \tableheader{Description} \\
\midrule
Lift & Shadows & Shifts the darkest values. Moving toward a color tints the shadows. \\
Gamma & Midtones & Adjusts the midrange brightness and color. The workhorse of color grading. \\
Gain & Highlights & Scales the brightest values. Moving toward a color tints the highlights. \\
\bottomrule
\end{tabularx}
\vspace{0.5em}

The node offers two mathematical modes:

\begin{itemize}[leftmargin=2em]
  \item \textbf{Lift/Gamma/Gain} --- the ASC CDL (American Society of Cinematographers Color Decision List) standard, using additive lift.
  \item \textbf{Offset/Power/Slope} --- an alternative color science model that some colorists prefer for its more intuitive behavior in certain correction scenarios.
\end{itemize}

\subsubsection{RGB Curves}

The \textbf{RGB Curves}\index{RGB Curves node} node provides per-channel curve adjustment, analogous to the Curves dialog in Photoshop:

\begin{itemize}[leftmargin=2em]
  \item \textbf{C (Combined)} --- adjusts all channels simultaneously for brightness and contrast control.
  \item \textbf{R, G, B} --- adjusts individual color channels for targeted color grading and cross-processing effects.
\end{itemize}

Common techniques:
\begin{itemize}[leftmargin=2em]
  \item Apply an S-curve on the Combined channel for a contrast boost.
  \item Lift the R channel shadows for warm shadow tones.
  \item Lower the B channel highlights for amber/warm highlights.
  \item Lift the bottom of all channels for a ``lifted blacks'' vintage film look.
\end{itemize}

\subsubsection{Hue/Saturation/Value}

The \textbf{Hue Saturation Value}\index{HSV node} node adjusts the HSV components of the image:

\vspace{0.5em}
\rowcolors{2}{rowgray}{white}
\begin{tabularx}{\textwidth}{l l X}
\toprule
\rowcolor{primary}
\tableheader{Control} & \tableheader{Range} & \tableheader{Effect} \\
\midrule
Hue & 0.0--1.0 (0.5 = no change) & Shifts all hues around the color wheel \\
Saturation & 0.0--2.0 (1.0 = no change) & Decreases or increases color intensity \\
Value & 0.0--2.0 (1.0 = no change) & Darkens or brightens the image \\
\bottomrule
\end{tabularx}
\vspace{0.5em}

\subsection{Effect Nodes}

\subsubsection{Glare}
\label{subsec:glare-node}

The \textbf{Glare}\index{Glare node} node creates light bloom, streaks, and glow effects from bright areas of the image. It is one of the most visually impactful compositor nodes and essential for selling the illusion of bright light sources:

\vspace{0.5em}
\rowcolors{2}{rowgray}{white}
\begin{tabularx}{\textwidth}{l L{3cm} X}
\toprule
\rowcolor{primary}
\tableheader{Mode} & \tableheader{Effect} & \tableheader{Use Case} \\
\midrule
Bloom & Soft glow around bright areas & General light glow, neon signs, magic effects \\
Streaks & Directional light streaks radiating from bright points & Lens artifacts, star filters, anamorphic streaks \\
Fog Glow & Soft, wide atmospheric glow & Dreamy atmosphere, fog, mist effects \\
Ghost & Concentric ring artifacts around bright points & Lens flare simulation \\
\bottomrule
\end{tabularx}
\vspace{0.5em}

Key parameters:
\begin{itemize}[leftmargin=2em]
  \item \textbf{Threshold} --- minimum brightness to trigger glare. Lower values cause more areas to glow; higher values restrict the effect to only the brightest pixels.
  \item \textbf{Quality} --- computation quality (High/Medium/Low). Use Low for previewing, High for final output.
  \item \textbf{Mix} --- blend amount with original image ($-1$ = glare only, $0$ = additive blend, $1$ = original only).
\end{itemize}

\begin{furrybox}
\textbf{Glare for Magical and Bioluminescent Fur Effects}

The Glare node in Bloom or Fog Glow mode is essential for selling bioluminescent or magical fur effects. Set your fur material's emission value high enough to exceed the Glare threshold, then use the Bloom mode with a moderate threshold (0.5--1.0) for a soft ethereal glow that wraps around individual fur strands. For convention badge art with glowing elements, the Fog Glow mode creates a wider, dreamier halo that works well at badge print resolution. See Chapter~\ref{ch:furryarts} for fur material setup.
\end{furrybox}

\subsubsection{Lens Distortion}

The \textbf{Lens Distortion}\index{Lens Distortion node} node simulates camera lens imperfections --- critical for matching CG renders to real camera footage:

\vspace{0.5em}
\rowcolors{2}{rowgray}{white}
\begin{tabularx}{\textwidth}{l X}
\toprule
\rowcolor{primary}
\tableheader{Parameter} & \tableheader{Effect} \\
\midrule
Distortion & Barrel (positive) or pincushion (negative) distortion \\
Dispersion & Chromatic aberration --- separates color channels at edges for a photographic look \\
Jitter & Adds noise to the distortion for more organic, less mathematically perfect results \\
Projector & Simulates projector-style distortion \\
\bottomrule
\end{tabularx}
\vspace{0.5em}

\begin{tipbox}
Adding a small amount of Lens Distortion (Distortion: 0.01--0.03, Dispersion: 0.005--0.01) to a perfectly clean CG render makes it feel more ``photographed'' and less ``computed.'' This subtle imperfection triggers the viewer's subconscious expectation of camera artifacts and increases the perceived realism of the image.
\end{tipbox}

\subsubsection{Defocus}

The \textbf{Defocus}\index{Defocus node} node creates depth-of-field blur in compositing using the Z-depth pass. While Cycles can render DOF natively via Camera Properties, compositing-based DOF offers the advantage of post-render control --- you can change the focal distance and aperture without re-rendering:

\vspace{0.5em}
\rowcolors{2}{rowgray}{white}
\begin{tabularx}{\textwidth}{l X}
\toprule
\rowcolor{primary}
\tableheader{Parameter} & \tableheader{Effect} \\
\midrule
Use Z Buffer & Whether to use the depth pass for variable blur across the image \\
Z Scale & Scale factor for depth values (adjusts the effective depth range) \\
f-Stop & Virtual aperture size (lower = shallower depth of field, more blur) \\
Max Blur & Maximum blur radius in pixels (limits computation time) \\
Threshold & Minimum difference in Z to apply blur (prevents halo artifacts at depth edges) \\
Bokeh Type & Shape of the blur kernel: Disk, Triangle, Square, Pentagon, Hexagon, Heptagon, Octagon \\
\bottomrule
\end{tabularx}
\vspace{0.5em}

The Bokeh Type parameter deserves special attention: it controls the shape of out-of-focus highlights. Real camera lenses produce bokeh shapes determined by the number of aperture blades. \textbf{Disk} produces smooth circular bokeh (simulating modern lenses with many blades), \textbf{Hexagon} creates the classic 6-blade aperture look of older lenses, and \textbf{Pentagon} gives a 5-blade character. Choose based on the aesthetic you want to match.

\subsection{Masking and Compositing Nodes}

\subsubsection{Cryptomatte Node}
\label{subsec:cryptomatte-node}

The \textbf{Cryptomatte}\index{Cryptomatte node} node extracts mattes from Cryptomatte render passes without requiring pre-render object index setup:

\begin{enumerate}[leftmargin=2em]
  \item Add a Cryptomatte node (Shift+A $\rightarrow$ Matte $\rightarrow$ Cryptomatte)
  \item Connect the Render Layers node's Image and Cryptomatte passes to the node
  \item Use the \textbf{eyedropper} tool in the Cryptomatte node header to click objects in the backdrop
  \item Each clicked object is added to the matte. Click again to remove it.
  \item Outputs: \textbf{Image} (original), \textbf{Matte} (alpha mask), \textbf{Pick} (ID visualization)
\end{enumerate}

Advantages over the legacy Object/Material Index system:
\begin{itemize}[leftmargin=2em]
  \item Anti-aliased edges (no staircase artifacts)
  \item Motion blur and depth of field preserved in the matte
  \item Transparency and overlapping objects handled correctly
  \item No per-object setup needed before rendering --- just enable the Cryptomatte pass and pick objects in compositing
\end{itemize}

\subsubsection{Alpha Over}

The \textbf{Alpha Over}\index{Alpha Over node} node composites one image over another using alpha transparency --- the fundamental operation of layered compositing:

\vspace{0.5em}
\rowcolors{2}{rowgray}{white}
\begin{tabularx}{\textwidth}{l X}
\toprule
\rowcolor{primary}
\tableheader{Input} & \tableheader{Purpose} \\
\midrule
Factor & Overall blend strength (0 = show bottom only, 1 = show top with alpha) \\
Image (bottom) & Background image \\
Image (top) & Foreground image (with alpha channel) \\
\bottomrule
\end{tabularx}
\vspace{0.5em}

\begin{warningbox}
\textbf{Premultiplied vs.\ Straight Alpha:} Enable the \textbf{Convert Premul} option when the top image uses premultiplied alpha (which is the default for most Blender render output). Leave it disabled for straight alpha images. Getting this wrong produces visible dark fringing around composited edges --- one of the most common compositing errors.
\end{warningbox}

\subsubsection{Mix (Blend Modes)}

The \textbf{Mix}\index{Mix node} node combines two images using standard blend modes. These are the same blend modes found in Photoshop, GIMP, and Krita:

\vspace{0.5em}
\rowcolors{2}{rowgray}{white}
\begin{tabularx}{\textwidth}{l X}
\toprule
\rowcolor{primary}
\tableheader{Blend Mode} & \tableheader{Effect} \\
\midrule
Mix & Simple linear interpolation between images \\
Add & Brightens by adding pixel values \\
Subtract & Darkens by subtracting pixel values \\
Multiply & Darkens; multiplies pixel values (ideal for shadows, AO overlays) \\
Screen & Brightens; inverse of multiply (ideal for glow, light effects) \\
Overlay & Contrast boost: multiply for darks, screen for lights \\
Soft Light & Gentle contrast and color tinting \\
Color & Applies hue and saturation from top image with luminance from bottom \\
Value & Applies luminance from top image with hue/saturation from bottom \\
\bottomrule
\end{tabularx}
\vspace{0.5em}

\subsubsection{Mask}

The \textbf{Mask}\index{Mask node} node loads a mask created in the Movie Clip Editor's mask tools. Masks are vector-based (spline) shapes that can be animated and feathered. They are used for rotoscoping, selective effects application, and region-based adjustments. Masks created for motion tracking purposes (see Section~\ref{sec:motion-tracking}) can also be referenced by this node.

\subsubsection{Viewer}

The \textbf{Viewer}\index{Viewer node} node is a critical workflow tool rather than an output node. Connect any node's output to a Viewer to see its result in the backdrop (if Backdrop is enabled) or in the Image Editor. The keyboard shortcut \textbf{Ctrl+Shift+Click} on any node automatically inserts a Viewer connection --- this is by far the fastest way to preview intermediate results while building complex node trees.

\subsection{Filter Nodes}

\vspace{0.5em}
\rowcolors{2}{rowgray}{white}
\begin{tabularx}{\textwidth}{l L{3.5cm} X}
\toprule
\rowcolor{primary}
\tableheader{Node} & \tableheader{Effect} & \tableheader{Use Case} \\
\midrule
Blur & Gaussian, Box, or other blur kernel & Softening, glow base, pre-processing for edge detection \\
Denoise & AI-based denoising (OptiX or OpenImageDenoise) & Removing path tracing noise from Cycles renders \\
Sharpen (via Filter) & Unsharp mask sharpening & Final output sharpening for web or print delivery \\
Dilate/Erode & Expands or contracts mask edges & Fixing edge fringing on mattes, growing or shrinking selections \\
Bilateral Blur & Edge-preserving blur & Smoothing skin or surfaces without losing hard edges \\
Pixelate & Reduces resolution & Retro/pixel art effects, stylistic choices \\
Sun Beams & Creates volumetric light rays from a source point & God rays, sunlight through windows, dramatic lighting \\
\bottomrule
\end{tabularx}
\vspace{0.5em}

%% ═══════════════════════════════════════════════════════════════════════════
\section{Practical Compositing Workflow}
\label{sec:practical-compositing}
%% ═══════════════════════════════════════════════════════════════════════════

\subsection{Standard Post-Production Pipeline}

The following pipeline represents the standard professional workflow for compositing CG renders in Blender:

\begin{codebox}
Render (Cycles / Eevee)\\
\hspace*{1em}|\\
\hspace*{1em}v\\
Save as Multi-Layer EXR\\
\hspace*{1em}|\\
\hspace*{1em}v\\
Load into Compositor (Image node)\\
\hspace*{1em}|\\
\hspace*{1em}v\\
Denoise (Denoise node using Denoising Data passes)\\
\hspace*{1em}|\\
\hspace*{1em}v\\
Recombine passes (if adjusted individually)\\
\hspace*{1em}|\\
\hspace*{1em}v\\
Color Grade (Color Balance, RGB Curves)\\
\hspace*{1em}|\\
\hspace*{1em}v\\
Effects (Glare, Lens Distortion, Vignette)\\
\hspace*{1em}|\\
\hspace*{1em}v\\
Output (Composite node -> final format)
\end{codebox}

\begin{tipbox}
Always save your renders as Multi-Layer EXR \textit{before} compositing. This separates the expensive render step from the iterative compositing step. You render once, then adjust compositing as many times as you want without re-rendering. This is not optional in professional work --- it is the standard.
\end{tipbox}

\subsection{Denoising Workflow}

\begin{enumerate}[leftmargin=2em]
  \item Enable the \textbf{Denoising Data} pass in View Layer Properties $\rightarrow$ Passes
  \item Render and save as Multi-Layer EXR
  \item In the Compositor, add a \textbf{Denoise} node
  \item Connect the Render Layers outputs:
  \begin{itemize}[leftmargin=2em]
    \item \textbf{Image} $\rightarrow$ Denoise \textbf{Image} input
    \item \textbf{Denoising Normal} $\rightarrow$ Denoise \textbf{Normal} input
    \item \textbf{Denoising Albedo} $\rightarrow$ Denoise \textbf{Albedo} input
  \end{itemize}
  \item Connect the Denoise \textbf{Image} output to the rest of the node tree
\end{enumerate}

The Denoise node uses Intel Open Image Denoise (OIDN) or NVIDIA OptiX for AI-based denoising that preserves detail while removing noise. The auxiliary passes (Normal and Albedo) provide the denoiser with geometric context, allowing it to distinguish between actual surface detail and noise --- this is why the Denoising Data pass must be enabled before rendering.

\subsection{Pass Recombination Workflow}

When you need to adjust individual light passes after rendering:

\begin{enumerate}[leftmargin=2em]
  \item Separate passes using the Render Layers node outputs
  \item Adjust individual passes (e.g., multiply Glossy Light by 0.5 to reduce reflections)
  \item Recombine using \textbf{Mix (Multiply)} for each color $\times$ light pair
  \item Use \textbf{Mix (Add)} to sum all light contributions back together
\end{enumerate}

\begin{codebox}
Diffuse Light --(adjust)--> Mix Multiply <-- Diffuse Color\\
\hspace*{12em}|\\
Glossy Light --(adjust)-->  Mix Multiply <-- Glossy Color\\
\hspace*{12em}|\\
\hspace*{7em}Mix Add (sum all) --> Color Grade --> Composite
\end{codebox}

\subsection{Background Replacement Workflow}

Compositing a CG render over a photographic background:

\begin{enumerate}[leftmargin=2em]
  \item Render the scene with a transparent background (\textbf{Film} $\rightarrow$ \textbf{Transparent} in Render Properties)
  \item Load a background image via an Image node
  \item Use the \textbf{Alpha Over} node: background on bottom input, render on top input
  \item Add \textbf{Lens Distortion} to match the camera characteristics of the background photograph
  \item Add \textbf{Color Balance} to match color temperature between the CG elements and the background
\end{enumerate}

\begin{furrybox}
\textbf{Badge Art Post-Processing Pipeline}

Convention badge art and character reference sheets benefit enormously from compositor post-processing. Render your character on a transparent background, then composite over a custom background plate. Apply a subtle Glare (Bloom mode, high threshold) to catch-light areas for eye sparkle. Add a gentle vignette (see Section~\ref{subsec:vignette}) to draw focus to the character. Use Color Balance to match the character's warmth to the background. Finally, add a small amount of Lens Distortion (Dispersion: 0.005) for photographic authenticity. This pipeline takes 5 minutes to set up and transforms a raw render into a polished portfolio piece.
\end{furrybox}

%% ═══════════════════════════════════════════════════════════════════════════
\section{Color Grading Fundamentals}
\label{sec:color-grading}
%% ═══════════════════════════════════════════════════════════════════════════

Color grading is the process of adjusting the color and tone of an image to establish mood, direct attention, and create visual consistency. It is the difference between a technically correct render and an emotionally compelling one. The same scene can feel cold and clinical or warm and inviting based entirely on color grading decisions.

\subsection{The Three-Way Color Corrector}

The Color Balance node in Lift/Gamma/Gain mode (see Section~\ref{subsec:color-balance}) is the standard tool for color grading in Blender's compositor. The three-way corrector divides the tonal range into three zones and allows independent color and brightness adjustments in each:

\begin{itemize}[leftmargin=2em]
  \item \textbf{Shadows (Lift):} Controls the darkest areas. Pulling toward blue creates ``cinematic'' cool shadows. Pulling toward orange creates warm, nostalgic shadows.
  \item \textbf{Midtones (Gamma):} The most visible adjustment --- it affects the overall color cast of the image. Subtle shifts here change the entire mood of the scene.
  \item \textbf{Highlights (Gain):} Controls bright areas. Warm highlights combined with cool shadows is the foundation of the ``teal and orange'' look that dominates modern cinema.
\end{itemize}

\subsection{Common Color Grading Looks}

\vspace{0.5em}
\rowcolors{2}{rowgray}{white}
\begin{tabularx}{\textwidth}{l L{1.6cm} L{1.8cm} L{1.8cm} X}
\toprule
\rowcolor{primary}
\tableheader{Look} & \tableheader{Shadows} & \tableheader{Midtones} & \tableheader{Highlights} & \tableheader{Additional} \\
\midrule
Teal \& Orange & Teal/blue & Slight warm & Orange/warm & Increase contrast with RGB Curves \\
Bleach Bypass & Crushed blacks & Desaturated & Bright whites & Reduce saturation to 0.5, increase contrast \\
Cross Process & Green/yellow & Shift hue & Magenta & Use RGB Curves per-channel \\
Day for Night & Deep blue & Strong blue & Blue-purple & Reduce exposure significantly, heavy blue tint \\
Vintage Film & Lifted blacks & Warm & Soft rolloff & Lift bottom of RGB Curves; reduce contrast \\
\bottomrule
\end{tabularx}
\vspace{0.5em}

\subsection{Scopes}

Blender provides scopes in the Image Editor and compositor sidebar for objective evaluation of color and exposure:

\vspace{0.5em}
\rowcolors{2}{rowgray}{white}
\begin{tabularx}{\textwidth}{l X}
\toprule
\rowcolor{primary}
\tableheader{Scope} & \tableheader{Purpose} \\
\midrule
Histogram & Shows value distribution per channel. Identifies clipping (values crushed to 0 or 1). \\
Waveform & Shows luminance or RGB values plotted against horizontal position. The industry standard for exposure evaluation. \\
Vectorscope & Shows hue and saturation distribution on a circular plot. Used for skin tone evaluation and overall color balance assessment. \\
\bottomrule
\end{tabularx}
\vspace{0.5em}

\begin{tipbox}
Never color grade by eye alone --- always reference the scopes. Your monitor's color accuracy, ambient lighting, and visual fatigue all affect your perception. The Waveform scope is your ground truth for exposure: if the waveform shows values hitting 0.0 (black) or 1.0 (white), you are clipping and losing data. The Vectorscope is essential for skin tones: regardless of ethnicity, all human skin tones fall on a specific line on the vectorscope (roughly between yellow and red). If your skin tones deviate from this line, your color grade is introducing an unnatural cast.
\end{tipbox}

%% ═══════════════════════════════════════════════════════════════════════════
\section{Video Sequence Editor (VSE)}
\label{sec:vse}
%% ═══════════════════════════════════════════════════════════════════════════

\subsection{Overview}

The \textbf{Video Sequence Editor (VSE)}\index{Video Sequence Editor}\index{VSE} is Blender's built-in non-linear video editor. It operates on \textbf{strips} (clips) arranged on \textbf{channels} (tracks). While not as feature-rich as dedicated NLEs like DaVinci Resolve or Premiere Pro, the VSE handles basic to intermediate editing tasks entirely within Blender --- which is invaluable when your 3D renders, compositing, and final edit all need to live in one pipeline.

Access the VSE via the \textbf{Video Editing} workspace or by changing an editor area to the Sequencer type.

\subsection{VSE Layout}

The Video Editing workspace typically shows three panels:

\vspace{0.5em}
\rowcolors{2}{rowgray}{white}
\begin{tabularx}{\textwidth}{l X}
\toprule
\rowcolor{primary}
\tableheader{Panel} & \tableheader{Purpose} \\
\midrule
Sequencer & Timeline view with strips arranged on channels \\
Preview & Visual preview of the current frame output \\
Properties (N-panel) & Strip properties, proxy settings, modifiers \\
\bottomrule
\end{tabularx}
\vspace{0.5em}

\subsection{Strip Types}

\vspace{0.5em}
\rowcolors{2}{rowgray}{white}
\begin{tabularx}{\textwidth}{l X l}
\toprule
\rowcolor{primary}
\tableheader{Strip Type} & \tableheader{Description} & \tableheader{Added Via} \\
\midrule
Movie & Video file (MP4, AVI, MOV, MKV, etc.) & Add $\rightarrow$ Movie \\
Image/Sequence & Single image or numbered image sequence & Add $\rightarrow$ Image/Sequence \\
Sound & Audio file (WAV, MP3, FLAC, OGG, etc.) & Add $\rightarrow$ Sound \\
Scene & Renders a Blender scene inline (live 3D render) & Add $\rightarrow$ Scene \\
Color & Solid color strip (backgrounds, fades) & Add $\rightarrow$ Color Strip \\
Text & Text overlay with font, size, color, shadow & Add $\rightarrow$ Text \\
Adjustment Layer & Applies effects to all strips below it on the timeline & Add $\rightarrow$ Adjustment Layer \\
Effect Strip & Various effects (Cross, Wipe, Transform, Speed, Glow, etc.) & Add $\rightarrow$ Effect Strip \\
\bottomrule
\end{tabularx}
\vspace{0.5em}

The \textbf{Scene} strip type deserves special mention: it renders a Blender scene directly into the VSE timeline, meaning you can intercut between live-rendered 3D scenes and imported video footage without any intermediate export step. This is a unique capability that dedicated NLEs cannot match.

\subsection{Essential VSE Keyboard Shortcuts}

\begin{shortcutbox}
\textbf{VSE Keyboard Shortcuts}

\vspace{0.3em}
\rowcolors{2}{rowgray}{white}
\begin{tabularx}{\textwidth}{l X}
\toprule
\rowcolor{primary}
\tableheader{Shortcut} & \tableheader{Action} \\
\midrule
Shift+A & Add Strip menu \\
G & Grab (move) selected strips along the timeline \\
K & Soft Cut (split strip at playhead, creates two linked strips) \\
Shift+K & Hard Cut (split strip at playhead, creates two independent strips) \\
X / Delete & Delete selected strips \\
Shift+D & Duplicate strips \\
E & Extend strip (change end frame) \\
H & Mute selected strips \\
Alt+H & Unmute all strips \\
Ctrl+C / Ctrl+V & Copy / Paste strips \\
Home & View All (zoom to fit all strips) \\
Page Up / Page Down & Move strip to higher/lower channel \\
\bottomrule
\end{tabularx}
\end{shortcutbox}

%% ═══════════════════════════════════════════════════════════════════════════
\section{VSE Editing Techniques}
\label{sec:vse-techniques}
%% ═══════════════════════════════════════════════════════════════════════════

\subsection{Cutting and Trimming}

\textbf{Soft Cut (K):} Splits a strip at the playhead into two strips that both reference the same source media. The first strip's end and the second strip's start are at the cut point. This is non-destructive --- the original media is unchanged.

\textbf{Hard Cut (Shift+K):} Like Soft Cut, but creates fully independent strips with no link to each other.

\textbf{Trimming:} Drag the left or right handle of a strip to reveal more or less of the source media:

\begin{itemize}[leftmargin=2em]
  \item Dragging the \textbf{left handle right} trims the beginning (the clip starts later in the source media).
  \item Dragging the \textbf{right handle left} trims the end (the clip ends earlier in the source media).
  \item Hold \textbf{Ctrl} while dragging to snap to the playhead or other strip edges.
\end{itemize}

\subsection{Transitions}

\subsubsection{Cross Dissolve}

\begin{enumerate}[leftmargin=2em]
  \item Overlap two strips on adjacent channels (the strips must overlap in time)
  \item Select both strips
  \item Add $\rightarrow$ Effect Strip $\rightarrow$ Cross
  \item The Cross strip automatically spans the overlap region
  \item The overlap duration controls the transition length
\end{enumerate}

\subsubsection{Wipe}

\begin{enumerate}[leftmargin=2em]
  \item Same setup as Cross Dissolve (overlapping strips on adjacent channels)
  \item Add $\rightarrow$ Effect Strip $\rightarrow$ Wipe
  \item Configure the wipe type: Single, Double, Iris, Clock
  \item Set direction and blur amount in the strip properties
\end{enumerate}

\subsubsection{Gamma Cross}

Like Cross Dissolve, but applies gamma correction during the dissolve for perceptually correct brightness transitions. A standard linear cross dissolve causes a visible ``dark dip'' in the middle of the transition because linear blending of gamma-encoded values produces darker-than-expected results. Gamma Cross corrects for this.

\begin{tipbox}
Always prefer \textbf{Gamma Cross} over regular Cross for dissolves unless you have a specific reason to use linear blending. The dark dip from linear cross dissolves is one of those subtle visual artifacts that audiences feel even if they cannot name it.
\end{tipbox}

\subsection{Speed Effects and Retiming}

\subsubsection{Speed Control (Legacy)}

The \textbf{Speed Control} effect strip modifies playback speed:

\begin{enumerate}[leftmargin=2em]
  \item Select a movie strip
  \item Add $\rightarrow$ Effect Strip $\rightarrow$ Speed Control
  \item Set speed factor: 2.0 = double speed, 0.5 = half speed
  \item For smooth slow motion, enable \textbf{Interpolation} (blends between source frames)
\end{enumerate}

\subsubsection{Retiming (Blender 4.0+)}

Blender 4.0 introduced native \textbf{retiming}\index{retiming} on strips, providing much more flexible speed control than the legacy Speed Control effect:

\begin{enumerate}[leftmargin=2em]
  \item Select a strip
  \item Open the Retiming panel in strip properties
  \item Add retiming keys at different points along the strip
  \item Set different speeds between keys
  \item Add speed transitions for smooth acceleration/deceleration between speed segments
\end{enumerate}

This system allows variable-speed effects within a single strip --- slow motion for an impact moment, normal speed for dialogue, fast motion for a montage --- all without splitting the strip or stacking effect layers.

\begin{furrybox}
\textbf{Furry Music Video Editing Pipeline}

The VSE is well-suited for editing furry music videos (FMVs) and dance compilation reels --- a popular content format in the furry community. Import your music as a Sound strip on channel 1, enable waveform display for beat-matching, then cut your video clips to the beat using K at each downbeat. Use the Retiming system (Blender 4.0+) for speed ramps that emphasize dance moves --- slow to 0.5x for acrobatic moments, then snap back to 1.0x on the beat. Layer Adjustment strips with Color Balance modifiers on top to shift the color mood per song section. The Scene strip type lets you intercut live 3D renders of your character model with the recorded footage for hybrid 2D/3D music videos.
\end{furrybox}

\subsection{Color Grading in VSE}

Each strip in the VSE can have \textbf{Modifiers} applied --- these work like compositor nodes but are attached directly to individual strips:

\vspace{0.5em}
\rowcolors{2}{rowgray}{white}
\begin{tabularx}{\textwidth}{l X}
\toprule
\rowcolor{primary}
\tableheader{Modifier} & \tableheader{Effect} \\
\midrule
Color Balance & Lift/Gamma/Gain three-way color correction \\
Curves & Per-channel curve adjustment \\
Hue Correct & Adjusts properties (saturation, value) based on input hue \\
Bright/Contrast & Simple brightness and contrast adjustment \\
White Balance & Adjust white point temperature \\
Tone Map & HDR to SDR tone mapping \\
\bottomrule
\end{tabularx}
\vspace{0.5em}

Modifiers stack and process in order from top to bottom. Use \textbf{Adjustment Layer} strips to apply modifiers to multiple strips simultaneously --- place the Adjustment Layer on a higher channel, spanning the range of strips you want to affect, and any modifiers on the Adjustment Layer will be applied to all visible strips below it.

\subsection{Text Strips}

Text strips create title cards, lower thirds, and text overlays:

\vspace{0.5em}
\rowcolors{2}{rowgray}{white}
\begin{tabularx}{\textwidth}{l X}
\toprule
\rowcolor{primary}
\tableheader{Property} & \tableheader{Description} \\
\midrule
Text & The text content (supports line breaks) \\
Font & Font family selection from installed system fonts \\
Size & Font size in pixels \\
Color & Text color \\
Shadow & Drop shadow with configurable offset and color \\
Box & Background box behind text (for readability over video) \\
Location & X/Y position (0.5, 0.5 = centered) \\
Alignment & Left, Center, Right \\
\bottomrule
\end{tabularx}
\vspace{0.5em}

Text properties can be keyframed for animated titles, enabling basic motion graphics directly within the VSE.

%% ═══════════════════════════════════════════════════════════════════════════
\section{H.265/HEVC Codec Support}
\label{sec:hevc}
%% ═══════════════════════════════════════════════════════════════════════════

\subsection{New in Blender 4.4}

Blender 4.4 added native support for \textbf{H.265/HEVC}\index{H.265}\index{HEVC} (High Efficiency Video Coding) as an output codec. This was a long-requested feature that brings Blender's video encoding capabilities to parity with modern delivery requirements.

\subsection{Configuration}

In \textbf{Output Properties} $\rightarrow$ \textbf{Output}:

\begin{itemize}[leftmargin=2em]
  \item Set Container to \textbf{MPEG-4} or \textbf{Matroska}
  \item Set Video Codec to \textbf{H.265}
  \item Configure quality settings:
  \begin{itemize}[leftmargin=2em]
    \item \textbf{Constant Rate Factor (CRF):} Lower = higher quality, larger file. Values of 18--23 are visually lossless for most content. CRF 0 is mathematically lossless.
    \item \textbf{Encoding Speed:} Slower presets produce smaller files at equivalent quality but take significantly longer to encode. Use ``Medium'' for general work, ``Slow'' or ``Slower'' for final delivery.
  \end{itemize}
\end{itemize}

\subsection{Bit Depth Support}

Blender 4.4 significantly expanded bit depth support across codecs:

\vspace{0.5em}
\rowcolors{2}{rowgray}{white}
\begin{tabularx}{\textwidth}{l C{2cm} C{2cm} C{2cm}}
\toprule
\rowcolor{primary}
\tableheader{Codec} & \tableheader{8-bit} & \tableheader{10-bit} & \tableheader{12-bit} \\
\midrule
H.264 & Yes & Yes (4.4+) & No \\
H.265/HEVC & Yes & Yes (4.4+) & Yes (4.4+) \\
AV1 & Yes & Yes (4.4+) & Yes (4.4+) \\
\bottomrule
\end{tabularx}
\vspace{0.5em}

10-bit encoding reduces banding artifacts in gradients and provides better quality for professional workflows. This is particularly important for content with subtle color transitions --- skies, skin tones, and volumetric effects all benefit from the additional precision.

\subsection{Color Space}

Videos rendered from Blender 4.4 explicitly use the \textbf{BT.709} color space standard (used for HDTV and most web video), ensuring consistent color reproduction across players and platforms. This eliminates the ambiguity that previously existed when video players had to guess the color space of Blender's output.

\subsection{H.265 vs.\ H.264 Comparison}

\vspace{0.5em}
\rowcolors{2}{rowgray}{white}
\begin{tabularx}{\textwidth}{l L{4.5cm} L{4.5cm}}
\toprule
\rowcolor{primary}
\tableheader{Property} & \tableheader{H.264 (AVC)} & \tableheader{H.265 (HEVC)} \\
\midrule
Compression efficiency & Baseline & 25--50\% better at same quality \\
File size & Larger & Smaller for equivalent quality \\
Encoding speed & Faster & Slower (more complex algorithms) \\
Hardware decode & Universal & Widespread (most devices since 2018) \\
License & Established patent pool & Complex patent pool (MPEG LA + others) \\
Max resolution & 4K (practical limit) & 8K+ \\
\bottomrule
\end{tabularx}
\vspace{0.5em}

\begin{tipbox}
For social media delivery (YouTube, Twitter/X, Instagram), H.264 remains the safest choice due to universal hardware decode support. For archival or high-quality distribution where file size matters, H.265 provides substantially better compression. For future-proofing, consider AV1, which is royalty-free and offers compression efficiency comparable to H.265.
\end{tipbox}

%% ═══════════════════════════════════════════════════════════════════════════
\section{Audio Synchronization}
\label{sec:audio-sync}
%% ═══════════════════════════════════════════════════════════════════════════

\subsection{Audio in the VSE}

Sound strips in the VSE support a range of audio formats and controls:

\begin{itemize}[leftmargin=2em]
  \item \textbf{Format support:} WAV, MP3, FLAC, OGG, and other FFmpeg-supported formats
  \item \textbf{Volume control} per strip (dB scale) for mixing
  \item \textbf{Pan} (stereo left/right positioning)
  \item \textbf{Waveform display} for visual reference and beat-matching
  \item \textbf{Scrubbing} --- hearing audio as you drag the playhead through the timeline
\end{itemize}

\subsection{Audio Sync Settings}

In the Timeline header or Playback menu, three sync modes are available:

\vspace{0.5em}
\rowcolors{2}{rowgray}{white}
\begin{tabularx}{\textwidth}{l X}
\toprule
\rowcolor{primary}
\tableheader{Setting} & \tableheader{Behavior} \\
\midrule
No Sync & Audio plays without frame synchronization. May drift over time. \\
Frame Dropping & Drops visual frames to keep audio in sync. Audio is the master clock. \\
AV-sync & Blender adjusts playback to maintain audio/video sync. Recommended for editing. \\
\bottomrule
\end{tabularx}
\vspace{0.5em}

\begin{warningbox}
Always use \textbf{AV-sync} or \textbf{Frame Dropping} mode when editing with audio. The ``No Sync'' mode allows audio and video to drift apart, producing an edit that looks synchronized during preview but is actually misaligned. Frame Dropping is preferred when smooth audio playback is essential (music videos, dialogue scenes); AV-sync is preferred when visual frame accuracy matters more than smooth playback.
\end{warningbox}

\subsection{Audio Workflow Tips}

\begin{itemize}[leftmargin=2em]
  \item \textbf{Separate audio and video} before applying speed effects to avoid desynchronization. Change the speed of the video strip independently from the audio strip.
  \item Use \textbf{Sound strip offset} to fine-tune sync (positive values delay audio start).
  \item For lip sync work, use \textbf{AV-sync} mode and scrub through the audio waveform to find phoneme boundaries.
  \item The \textbf{Mixdown} button (VSE header) renders all audio strips to a single file for final export or for importing into external audio software for mastering.
  \item Audio can be cached in memory for faster scrubbing: \textbf{User Preferences} $\rightarrow$ \textbf{System} $\rightarrow$ \textbf{Audio} settings.
\end{itemize}

%% ═══════════════════════════════════════════════════════════════════════════
\section{Motion Tracking}
\label{sec:motion-tracking}
%% ═══════════════════════════════════════════════════════════════════════════

\begin{historynote}
Motion tracking --- the process of extracting camera or object movement from video footage --- was once the exclusive domain of facilities with six-figure software budgets. Boujou (2D3), PFTrack, and SynthEyes were the standard tools, each costing thousands of dollars per seat. Blender's inclusion of a complete motion tracking system (introduced in Blender 2.61, 2011) was a landmark moment in democratizing VFX tools. The tracking algorithm is based on the libmv library, originally developed by Sergey Sharybin, who joined the Blender Foundation as a core developer.
\end{historynote}

\subsection{Overview}

Blender includes a complete \textbf{motion tracking}\index{motion tracking} system for reconstructing camera and object movement from video footage. This is essential for VFX work --- adding 3D elements to live-action footage so they appear to exist in the real world. Access the tracking workspace via the \textbf{Motion Tracking} workspace tab.

\subsection{Camera Tracking Workflow}

Camera tracking (also called \textit{matchmoving}) reconstructs the 3D camera path from a 2D video. The process follows five steps:

\textbf{Step 1: Import Footage}
\begin{itemize}[leftmargin=2em]
  \item Open the Movie Clip Editor
  \item Open your footage file
  \item Set the camera intrinsics (focal length, sensor size) if known from the original camera --- this dramatically improves solve accuracy
\end{itemize}

\textbf{Step 2: Set Tracking Parameters}
\begin{itemize}[leftmargin=2em]
  \item In the Track panel, set \textbf{Pattern Size} --- the area around each tracking point that is matched frame-to-frame
  \item Set \textbf{Search Size} --- the area in the next frame where Blender looks for the pattern. Larger search areas handle faster camera movement but are slower to compute.
\end{itemize}

\textbf{Step 3: Place and Track Markers}
\begin{itemize}[leftmargin=2em]
  \item Place tracking markers on high-contrast features: corners, edges, distinct patterns
  \item Minimum 8 markers required, but 15--20 is recommended for robust solves
  \item Track forward/backward using the Track buttons or keyboard shortcuts
\end{itemize}

\begin{shortcutbox}
\textbf{Motion Tracking Shortcuts}

\vspace{0.3em}
\rowcolors{2}{rowgray}{white}
\begin{tabularx}{\textwidth}{l X}
\toprule
\rowcolor{primary}
\tableheader{Shortcut} & \tableheader{Action} \\
\midrule
Ctrl+T & Track selected markers forward \\
Shift+Ctrl+T & Track selected markers backward \\
Alt+T & Track all markers forward \\
E & Detect features (auto-place markers on high-contrast points) \\
\bottomrule
\end{tabularx}
\end{shortcutbox}

\textbf{Step 4: Solve Camera Motion}
\begin{itemize}[leftmargin=2em]
  \item Set two \textbf{Keyframes} --- frames with good marker distribution and significant camera movement between them
  \item Click \textbf{Solve Camera Motion}
  \item Check the \textbf{Solve Error} (reprojection error):
  \begin{itemize}[leftmargin=2em]
    \item Below 0.3 = excellent solve
    \item 0.3 to 3.0 = usable, may need marker refinement
    \item Above 3.0 = poor; remove bad markers and re-solve
  \end{itemize}
\end{itemize}

\textbf{Step 5: Set Up Scene}
\begin{itemize}[leftmargin=2em]
  \item Click \textbf{Set Up Tracking Scene} to automatically create a 3D scene with the solved camera and a ground plane
  \item Or manually: use \textbf{Set as Background} and \textbf{Set Origin} to align the 3D scene with the tracked footage
\end{itemize}

\subsection{Object Tracking}

Object tracking follows the motion of a specific rigid object within the footage, separate from the camera motion:

\begin{enumerate}[leftmargin=2em]
  \item Create a new \textbf{Object} in the tracking settings
  \item Place markers on the object (minimum 8 markers on the object itself)
  \item Track the markers through the footage
  \item Solve the object's motion separately from the camera solve
  \item The solved object appears as an Empty in the 3D scene, which can be used as a parent for 3D geometry
\end{enumerate}

\subsection{Plane Tracking}

\textbf{Plane tracking}\index{plane tracking} follows the motion of a flat surface --- a wall, floor, screen, or sign --- for projecting replacement content:

\begin{enumerate}[leftmargin=2em]
  \item Track at least 4 markers on the flat surface
  \item Select those markers
  \item Create a \textbf{Plane Track} (Header $\rightarrow$ Track $\rightarrow$ Plane Track $\rightarrow$ Create)
  \item The plane track can be used with the \textbf{Plane Track Deform} compositor node to project images onto the tracked surface
\end{enumerate}

Plane tracking is used for replacing screens, signs, billboards, and any other flat surface in footage with CG content.

\subsection{Masking}

The Movie Clip Editor includes a \textbf{Mask} tool for rotoscoping:

\begin{itemize}[leftmargin=2em]
  \item Create spline-based masks to isolate objects in footage
  \item Animate the mask shape per-frame for moving objects
  \item Masks can be used in the compositor via the \textbf{Mask} node (see Section~\ref{sec:compositor-nodes})
  \item Feathering controls create soft mask edges for natural-looking composites
\end{itemize}

%% ═══════════════════════════════════════════════════════════════════════════
\section{VFX Compositing with Tracked Footage}
\label{sec:vfx-compositing}
%% ═══════════════════════════════════════════════════════════════════════════

\subsection{Standard VFX Pipeline in Blender}

The following pipeline shows the complete VFX workflow from footage import to final output:

\begin{codebox}
Footage Import\\
\hspace*{1em}|\\
\hspace*{1em}v\\
Motion Tracking (Movie Clip Editor)\\
\hspace*{1em}|\\
\hspace*{1em}v\\
Camera Solve (produces solved camera)\\
\hspace*{1em}|\\
\hspace*{1em}v\\
Scene Setup (3D scene aligned to footage)\\
\hspace*{1em}|\\
\hspace*{1em}v\\
3D Modeling \& Animation (in the tracked scene)\\
\hspace*{1em}|\\
\hspace*{1em}v\\
Shadow Catchers \& Holdout Objects\\
\hspace*{1em}|\\
\hspace*{1em}v\\
Render (with transparent Film)\\
\hspace*{1em}|\\
\hspace*{1em}v\\
Composite (Compositor: CG over footage)\\
\hspace*{1em}|\\
\hspace*{1em}v\\
Final Output
\end{codebox}

\subsection{Shadow Catchers}

A \textbf{Shadow Catcher}\index{Shadow Catcher} is an object that receives shadows from CG objects but is itself invisible in the final render. This is the key technique for casting CG shadows onto live-action surfaces:

\begin{enumerate}[leftmargin=2em]
  \item Add a plane positioned where the ground surface is in the footage
  \item In Object Properties $\rightarrow$ Visibility, enable \textbf{Shadow Catcher}
  \item Render with Film $\rightarrow$ Transparent enabled
  \item The Shadow pass captures shadows cast onto the catcher
  \item In the compositor, multiply the shadow pass with the footage to overlay the CG shadows
\end{enumerate}

\subsection{Holdout Objects}

\textbf{Holdout}\index{Holdout} objects are invisible but block CG objects behind them. They are used when a real-world object needs to occlude (appear in front of) CG elements:

\begin{enumerate}[leftmargin=2em]
  \item Create a rough 3D model matching the real object's shape and position in the footage
  \item In Material Properties, set the surface shader to \textbf{Holdout}
  \item The holdout object cuts a hole in the CG render, letting the original footage show through at that location
\end{enumerate}

\begin{tipbox}
Holdout objects do not need to be detailed --- they just need to match the silhouette of the real-world object from the camera's perspective. A rough box or cylinder is often sufficient. Spend time on positioning accuracy, not on modeling detail.
\end{tipbox}

\subsection{Compositing CG with Footage}

The final compositing stage integrates CG renders with the original footage in the Compositor:

\begin{enumerate}[leftmargin=2em]
  \item Add a \textbf{Movie Clip} node (the original footage)
  \item Add the \textbf{Render Layers} node (the CG render with transparent background)
  \item Use \textbf{Alpha Over} to composite CG over footage
  \item Add \textbf{Color Balance} to match the CG color temperature to the footage
  \item Add \textbf{Lens Distortion} to match the CG to the footage camera's distortion characteristics
  \item Use the \textbf{Shadow} pass with a Mix (Multiply) node to overlay CG shadows onto the footage
  \item Add \textbf{Glare} sparingly for light integration (matching the bloom character of the footage's camera)
  \item Route the final result to the \textbf{Composite} node for output
\end{enumerate}

\begin{warningbox}
The most common VFX integration failure is \textbf{color mismatch} between CG and footage. Real cameras produce images with specific color characteristics (white balance, contrast curve, color science) that vary between camera manufacturers. If your CG render is perfectly neutral but the footage has a warm cast from tungsten lighting, the CG elements will look ``pasted on.'' Always use the Color Balance node to match the CG's color temperature, contrast, and saturation to the footage.
\end{warningbox}

%% ═══════════════════════════════════════════════════════════════════════════
\section{Grease Pencil}
\label{sec:grease-pencil}
%% ═══════════════════════════════════════════════════════════════════════════

\subsection{Overview}

\textbf{Grease Pencil}\index{Grease Pencil} is Blender's 2D drawing and animation system that operates within 3D space. Each stroke has XYZ position data, allowing 2D art to interact with 3D scenes --- receiving lighting, casting shadows, being occluded by geometry, and participating in camera movement. This hybrid 2D/3D approach is unique among major creative applications and enables workflows that are impossible in dedicated 2D animation software like Toon Boom or TVPaint.

\subsection{Grease Pencil 3.0 (Blender 4.3+)}

Blender 4.3 introduced \textbf{Grease Pencil 3.0}\index{Grease Pencil 3.0}, a complete rewrite with major performance and integration improvements:

\vspace{0.5em}
\rowcolors{2}{rowgray}{white}
\begin{tabularx}{\textwidth}{l L{4cm} L{4cm}}
\toprule
\rowcolor{primary}
\tableheader{Feature} & \tableheader{GP 2.x (Legacy)} & \tableheader{GP 3.0} \\
\midrule
Performance & Python-heavy, slow with many strokes & C++ core, significantly faster \\
Geometry Nodes & Limited support & Full Geometry Nodes integration \\
Data model & Custom data format & Standard Blender data (compatible with other systems) \\
Scalability & Performance degraded with density & Handles high-density drawings and long timelines \\
\bottomrule
\end{tabularx}
\vspace{0.5em}

The GP 3.0 rewrite is particularly significant because it brings Grease Pencil strokes into the same data model as meshes, curves, and other Blender objects. This means Geometry Nodes --- Blender's procedural modeling system --- can now operate on Grease Pencil data, opening up procedural 2D workflows that were previously impossible. See Chapter~\ref{ch:scripting} for Geometry Nodes details.

\subsection{Drawing Tools}

\vspace{0.5em}
\rowcolors{2}{rowgray}{white}
\begin{tabularx}{\textwidth}{l X}
\toprule
\rowcolor{primary}
\tableheader{Tool} & \tableheader{Purpose} \\
\midrule
Draw & Freehand stroke drawing with pressure sensitivity (requires tablet) \\
Fill & Fills enclosed areas with color (works like a paint bucket) \\
Erase & Removes strokes or portions of strokes \\
Tint & Applies color tinting to existing strokes without redrawing \\
Cutter & Cuts strokes at intersections with drawn lines \\
Line & Draws straight lines between two points \\
Arc & Draws arcs and curves with adjustable curvature \\
Box / Circle & Draws rectangles and circles (hold Shift for squares/perfect circles) \\
Polyline & Draws connected line segments (click for each vertex) \\
\bottomrule
\end{tabularx}
\vspace{0.5em}

\subsection{Stroke Sculpting}

In Sculpt Mode, Grease Pencil strokes can be deformed using specialized sculpt brushes, providing an organic approach to refining line work:

\vspace{0.5em}
\rowcolors{2}{rowgray}{white}
\begin{tabularx}{\textwidth}{l X}
\toprule
\rowcolor{primary}
\tableheader{Brush} & \tableheader{Effect} \\
\midrule
Smooth & Smooths stroke curves, removing jitter from freehand drawing \\
Thickness & Adjusts stroke width in the brush radius \\
Strength & Adjusts stroke opacity \\
Grab & Moves strokes by dragging, like pulling on a rubber sheet \\
Push & Pushes strokes in the brush direction \\
Twist & Rotates strokes around the brush center \\
Pinch & Pulls strokes toward the brush center \\
Randomize & Adds random displacement for a hand-drawn, organic feel \\
Clone & Stamps duplicates of selected strokes \\
\bottomrule
\end{tabularx}
\vspace{0.5em}

\subsection{Grease Pencil Materials}

Grease Pencil objects use their own material type with two components:

\vspace{0.5em}
\rowcolors{2}{rowgray}{white}
\begin{tabularx}{\textwidth}{l X}
\toprule
\rowcolor{primary}
\tableheader{Component} & \tableheader{Purpose} \\
\midrule
Stroke & Controls the outline/line appearance: color, opacity, line style \\
Fill & Controls the interior fill: solid color, gradient, texture, or checker pattern \\
\bottomrule
\end{tabularx}
\vspace{0.5em}

Materials can use \textbf{Holdout} mode to erase underlying strokes, creating complex layered effects --- for example, a ``window'' material that reveals a background layer through an upper drawing layer.

\subsection{Layers and Onion Skinning}

\textbf{Layers} organize strokes much like Photoshop layers. Each layer has independent opacity, blend mode, and lock settings. Use layers to separate line art, fills, effects, and backgrounds so they can be edited independently.

\textbf{Onion Skinning}\index{onion skinning} displays adjacent frames as transparent overlays for animation reference. This is indispensable for frame-by-frame animation, as it shows you where the drawing was on previous and upcoming frames:

\begin{itemize}[leftmargin=2em]
  \item Configurable colors for before/after frames (typically blue for past, red for future)
  \item Adjustable number of visible onion skin frames
  \item Per-frame opacity control
\end{itemize}

\subsection{Animation}

Grease Pencil animation is fundamentally frame-by-frame:

\begin{itemize}[leftmargin=2em]
  \item Each keyframe contains a \textbf{drawing} (a set of strokes)
  \item Blank keyframes create holds --- the current drawing persists until the next keyframe containing new strokes
  \item \textbf{Interpolation} between keyframes can be automatic (Blender attempts to morph between stroke shapes) or manual (the animator draws each frame)
\end{itemize}

For character animation, the combination of frame-by-frame drawing with onion skinning and GP 3.0's improved performance makes Blender a viable alternative to dedicated 2D animation software for many workflows. See Chapter~\ref{ch:animation} for animation principles that apply to both 3D and Grease Pencil animation.

%% ═══════════════════════════════════════════════════════════════════════════
\section{2D/3D Hybrid Animation Workflows}
\label{sec:hybrid-animation}
%% ═══════════════════════════════════════════════════════════════════════════

\subsection{Grease Pencil in Production}

Grease Pencil has proven itself in professional production. Most notably, \textit{Spider-Man: Across the Spider-Verse} (2023) used Blender's Grease Pencil for freehand drawing by animators to add 2D strokes and line effects to the 3D animation. These strokes were then converted to meshes and exported to the main Maya pipeline via custom scripts. The film's distinctive visual style --- where 2D line work, smears, and hand-drawn effects overlay 3D characters and environments --- demonstrated the production viability of the hybrid approach.

\subsection{Three Hybrid Workflow Approaches}

\subsubsection{Approach 1: 3D Base with 2D Overlay}

\begin{enumerate}[leftmargin=2em]
  \item Create the 3D scene with characters and environments using standard modeling and animation
  \item Add Grease Pencil objects for 2D effects: speed lines, impact frames, expression marks, action smears
  \item Because Grease Pencil strokes exist in 3D space, they move with the camera and maintain correct spatial relationships with the 3D scene
\end{enumerate}

This approach is ideal for adding hand-drawn expressiveness to otherwise photorealistic or stylized 3D animation.

\subsubsection{Approach 2: 2D Animation in 3D Environment}

\begin{enumerate}[leftmargin=2em]
  \item Model a 3D environment (sets, props, backgrounds)
  \item Create Grease Pencil characters animated frame-by-frame
  \item Characters benefit from 3D camera movement, lighting, and depth
  \item Use \textbf{Surface Offset} on Grease Pencil objects to prevent Z-fighting with 3D geometry
\end{enumerate}

This approach gives 2D characters the spatial richness of a 3D world --- camera moves that would require complex multiplane setups in traditional 2D animation come for free.

\subsubsection{Approach 3: Storyboard Animatic}

\begin{enumerate}[leftmargin=2em]
  \item Draw storyboard frames with Grease Pencil directly in the 3D viewport
  \item Time them on the Timeline to establish pacing
  \item Add 3D camera moves for cinematic framing
  \item Export as an animatic video for editorial review before committing to full production
\end{enumerate}

\begin{furrybox}
\textbf{Compositing Fur Renders with 2D Effects}

The hybrid 2D/3D workflow is particularly powerful for furry character animation. Render your 3D fur character (see Chapter~\ref{ch:shading} for hair/fur particle systems), then use Grease Pencil overlays for expressive 2D elements that are difficult to achieve in pure 3D: speed lines during fast movement, exaggerated expression marks (sweat drops, anger veins, sparkle effects), impact frames on key animation poses, and comic-style text bubbles. The Grease Pencil strokes exist in 3D space, so they correctly track with camera movement and maintain depth relative to the character. Use the Build modifier on the GP object to animate strokes appearing one at a time for dynamic effect reveals.
\end{furrybox}

\subsection{Grease Pencil Modifiers}

Grease Pencil objects support a specialized set of modifiers for non-destructive effects:

\vspace{0.5em}
\rowcolors{2}{rowgray}{white}
\begin{tabularx}{\textwidth}{l X}
\toprule
\rowcolor{primary}
\tableheader{Modifier} & \tableheader{Effect} \\
\midrule
Noise & Adds jitter to stroke position/thickness for a hand-drawn, sketchy feel \\
Smooth & Smooths stroke curves non-destructively \\
Thickness & Uniformly adjusts stroke width across all strokes \\
Tint & Applies a color overlay to all strokes \\
Color & Modifies hue, saturation, value of stroke colors \\
Opacity & Adjusts transparency of all strokes \\
Build & Progressively draws/reveals strokes over time (animated stroke reveal) \\
Mirror & Mirrors strokes across one or more axes \\
Array & Repeats strokes in patterns (grid, radial, etc.) \\
Offset & Translates strokes per-layer for parallax scrolling effects \\
Lattice & Deforms strokes using a Lattice object for broad shape changes \\
Armature & Deforms strokes using an Armature for skeletal animation of 2D art \\
\bottomrule
\end{tabularx}
\vspace{0.5em}

The \textbf{Armature} modifier is particularly noteworthy: it allows you to rig 2D Grease Pencil characters with the same skeletal animation system used for 3D characters. This means you can use Blender's full animation toolset --- inverse kinematics, constraints, drivers, the Graph Editor --- to animate 2D drawings.

\subsection{Shader Effects (Visual Effects for Grease Pencil)}

Grease Pencil objects support real-time visual effects applied at render time:

\vspace{0.5em}
\rowcolors{2}{rowgray}{white}
\begin{tabularx}{\textwidth}{l X}
\toprule
\rowcolor{primary}
\tableheader{Effect} & \tableheader{Description} \\
\midrule
Blur & Blurs the Grease Pencil render (depth-of-field simulation for 2D) \\
Colorize & Applies color filters: grayscale, sepia, duotone, and more \\
Flip & Mirrors the render horizontally or vertically \\
Glow & Adds outer glow to strokes (neon sign effect) \\
Light & Applies 2D lighting to Grease Pencil strokes \\
Pixelate & Reduces render resolution for pixel art aesthetic \\
Rim & Adds rim lighting effect to strokes \\
Shadow & Adds drop shadow to all strokes \\
Swirl & Applies swirl distortion from a center point \\
Wave & Applies wave distortion (animated water/heat haze effect) \\
\bottomrule
\end{tabularx}
\vspace{0.5em}

%% ═══════════════════════════════════════════════════════════════════════════
\section{Advanced Compositing Techniques}
\label{sec:advanced-compositing}
%% ═══════════════════════════════════════════════════════════════════════════

This section covers techniques that build on the fundamentals from earlier sections. Each technique is a practical recipe that you can implement directly.

\subsection{Fog and Atmospheric Depth}

Using the \textbf{Mist Pass} to create distance-based atmospheric effects:

\begin{enumerate}[leftmargin=2em]
  \item Enable \textbf{Mist Pass} in View Layer Properties $\rightarrow$ Passes
  \item Configure Mist settings in World Properties $\rightarrow$ Mist Pass:
  \begin{itemize}[leftmargin=2em]
    \item \textbf{Start} --- distance from camera where mist begins
    \item \textbf{Depth} --- distance over which mist fades from 0 to 1
    \item \textbf{Falloff} --- curve type: Quadratic (default, natural-looking), Linear, or Inverse Quadratic
  \end{itemize}
  \item In the Compositor, use the Mist pass as a Factor input on a \textbf{Mix} node
  \item Mix between the render and a fog color (typically a desaturated blue-gray)
\end{enumerate}

This creates realistic aerial perspective where distant objects fade into atmospheric haze. The technique is especially effective for landscape renders and urban scenes with depth, and is far cheaper to adjust in compositing than using volumetric fog in the scene itself.

\subsection{Depth of Field in Compositing}

While Cycles can render depth of field natively via Camera Properties, compositing-based DOF offers more control after rendering and avoids the render-time cost of in-camera DOF:

\begin{enumerate}[leftmargin=2em]
  \item Render with the \textbf{Z Depth} pass enabled
  \item Add a \textbf{Defocus} node in the Compositor
  \item Connect the Z pass to the Defocus node's Z input
  \item Set the \textbf{f-Stop} to control blur intensity
  \item Adjust \textbf{Max Blur} to prevent excessive computation
  \item Choose the \textbf{Bokeh Type} to match the desired lens character
  \item Set \textbf{Z Scale} to fine-tune the depth range that appears in focus
\end{enumerate}

\begin{tipbox}
For best results with compositing-based DOF, render at a higher resolution than your final output and use \textbf{Map Value} or \textbf{Math} nodes to remap the Z-depth range before feeding it to the Defocus node. The raw Z-depth values from Cycles can span a very large range (from near-zero to thousands of Blender units), and the Defocus node works best with a normalized depth range.
\end{tipbox}

\subsection{Edge Detection and Outline Rendering}

Creating NPR (Non-Photorealistic Rendering) outlines in the compositor --- useful for toon-shaded looks, technical illustration, and stylized renders:

\subsubsection{Method 1: Normal Pass Edge Detection}

\begin{enumerate}[leftmargin=2em]
  \item Enable the \textbf{Normal} pass
  \item Add a \textbf{Filter} node set to \textbf{Sobel} (edge detection algorithm)
  \item Connect the Normal pass to the Filter node
  \item The output highlights edges where surface normals change sharply
  \item Invert the result and multiply it over the render for a dark outline effect
\end{enumerate}

\subsubsection{Method 2: Depth-Based Outlines}

\begin{enumerate}[leftmargin=2em]
  \item Use the Z-Depth pass
  \item Apply a \textbf{Filter (Sobel)} node to detect depth discontinuities
  \item This creates outlines at silhouette edges where depth changes abruptly
  \item Combine with Normal-based edges for both silhouette and surface detail outlines
\end{enumerate}

\subsubsection{Method 3: Freestyle}

Blender's \textbf{Freestyle}\index{Freestyle} line renderer (Render Properties $\rightarrow$ Freestyle) generates vector-quality edge lines with extensive style control. While not a compositor technique per se, the resulting line render can be composited over the main render using Alpha Over for maximum control over line weight, color, and style independently from the 3D shading. See Chapter~\ref{ch:shading} for Freestyle configuration.

\subsection{Vignette Effect}
\label{subsec:vignette}

A vignette darkens the edges of the frame, drawing the viewer's eye toward the center. It is a simple but effective finishing touch:

\begin{enumerate}[leftmargin=2em]
  \item Add an \textbf{Ellipse Mask} node (Shift+A $\rightarrow$ Matte $\rightarrow$ Ellipse Mask)
  \item Set width and height to approximately 0.9 each
  \item Add a \textbf{Blur} node and connect the mask output to it
  \item Set Blur to Gaussian with a high radius (200+ pixels) for smooth falloff
  \item Use the blurred mask as a Factor input in a \textbf{Mix} node
  \item Mix between a darkened version of the image and the original
  \item The edges darken while the center remains bright
\end{enumerate}

\subsection{Light Wrap}

\textbf{Light Wrap}\index{light wrap} blends bright areas of a background plate into the edges of a composited foreground element, creating more natural integration. Without light wrap, composited elements can appear ``cut out'' and pasted onto the background:

\begin{enumerate}[leftmargin=2em]
  \item Blur the background plate heavily (Blur node, 50--100+ pixel radius)
  \item Use the foreground's alpha (dilated slightly with the \textbf{Dilate/Erode} node) as a mask
  \item Mix the blurred background into the foreground edges using \textbf{Alpha Over} or \textbf{Mix (Screen)}
  \item This simulates light bleeding from the background onto the foreground edges, matching how a camera would capture the scene
\end{enumerate}

\subsection{Multi-Pass Color Correction Workflow}

For maximum flexibility in professional work, correct color per-pass before recombining:

\vspace{0.5em}
\rowcolors{2}{rowgray}{white}
\begin{tabularx}{\textwidth}{l X}
\toprule
\rowcolor{primary}
\tableheader{Pass} & \tableheader{Typical Adjustments} \\
\midrule
Diffuse Light & Brighten fill light areas, cool shadow tones \\
Diffuse Color & Adjust skin tones, environment color matching \\
Glossy Light & Reduce or boost specular intensity; tint reflections \\
Emission & Bloom amount, glow color tinting \\
AO & Strengthen for stylistic darkness in crevices; weaken for flat/bright look \\
Shadow & Adjust shadow density for VFX matching \\
\bottomrule
\end{tabularx}
\vspace{0.5em}

This per-pass approach is standard in VFX pipelines and allows changes that would otherwise require a complete re-render.

\subsection{Node Groups and Reusability}

Create reusable compositing setups with \textbf{Node Groups}\index{Node Groups}:

\begin{enumerate}[leftmargin=2em]
  \item Select a group of connected nodes that form a logical unit
  \item Press \textbf{Ctrl+G} to create a Node Group
  \item The group becomes a single node with exposed inputs and outputs
  \item Name the group in the N-panel (e.g., ``Color Grade -- Warm'')
  \item Reuse the group in other compositing setups via Shift+A $\rightarrow$ Group
  \item Groups are stored as data blocks and can be shared between .blend files via File $\rightarrow$ Append
\end{enumerate}

Common reusable node groups to build for your personal library:

\begin{itemize}[leftmargin=2em]
  \item \textbf{Film Look} --- Curves + Color Balance + subtle Glare for a cinematic finish
  \item \textbf{Shadow Composite} --- standardized shadow pass integration with adjustable density
  \item \textbf{Matte Edge Clean} --- Dilate/Erode + Blur for consistent matte edge treatment
  \item \textbf{Vignette} --- parameterized vignette with adjustable size and falloff intensity
\end{itemize}

\subsection{Denoise Strategies}

\vspace{0.5em}
\rowcolors{2}{rowgray}{white}
\begin{tabularx}{\textwidth}{l X l}
\toprule
\rowcolor{primary}
\tableheader{Strategy} & \tableheader{When to Use} & \tableheader{Quality} \\
\midrule
Cycles Denoiser (viewport) & Quick preview during modeling/lighting & Low \\
Denoise node (OIDN) & General purpose post-render denoising & High \\
Denoise node (OptiX) & NVIDIA GPU acceleration available & High \\
Higher sample count & When denoising artifacts are unacceptable & Highest \\
Denoise per-pass & Individual passes denoised before recombination & Very High \\
\bottomrule
\end{tabularx}
\vspace{0.5em}

\begin{tipbox}
When denoising individual passes, apply the Denoise node to each light pass separately \textit{before} multiplying with the corresponding color pass. This prevents the denoiser from smearing fine color details that exist in the color pass. The light passes contain the noise (from insufficient path tracing samples), while the color passes are essentially noise-free --- denoising the combined result risks smoothing out the color detail along with the noise.
\end{tipbox}

%% ═══════════════════════════════════════════════════════════════════════════
\section{Common Pitfalls and Practical Tips}
\label{sec:compositing-pitfalls}
%% ═══════════════════════════════════════════════════════════════════════════

\subsection{Compositor Pitfalls}

\vspace{0.5em}
\rowcolors{2}{rowgray}{white}
\begin{tabularx}{\textwidth}{L{3cm} L{3.5cm} X}
\toprule
\rowcolor{primary}
\tableheader{Pitfall} & \tableheader{Cause} & \tableheader{Solution} \\
\midrule
Compositing not updating & Use Nodes not enabled, or node tree disconnected from Composite output & Check the Use Nodes checkbox; ensure the Composite node is connected to your node tree \\
Black output from Viewer & No Viewer node connected to the node you want to inspect & Ctrl+Shift+Click on a node to add a Viewer connection \\
EXR passes missing & Passes not enabled before rendering & Enable passes in View Layer Properties \textit{before} rendering \\
Alpha fringing & Straight/premultiplied alpha mismatch & Enable Convert Premul on the Alpha Over node; check image alpha type in Image Editor \\
Denoise node produces no change & Denoising Data passes not enabled & Enable Denoising Data in View Layer Properties $\rightarrow$ Passes before rendering \\
Glare too strong & Threshold too low & Increase Threshold; reduce Mix amount toward 1.0 \\
\bottomrule
\end{tabularx}
\vspace{0.5em}

\subsection{VSE Pitfalls}

\vspace{0.5em}
\rowcolors{2}{rowgray}{white}
\begin{tabularx}{\textwidth}{L{3cm} L{3.5cm} X}
\toprule
\rowcolor{primary}
\tableheader{Pitfall} & \tableheader{Cause} & \tableheader{Solution} \\
\midrule
Audio out of sync & Wrong sync mode selected & Set Sync Mode to AV-sync in the Timeline header \\
Video plays at wrong speed & Source FPS differs from project FPS & Match Output Properties $\rightarrow$ Frame Rate to the footage's frame rate \\
Quality loss on export & Encoding settings too aggressive & Lower the CRF value (higher quality); verify resolution matches \\
Strips overlapping incorrectly & Channel order confusion & Higher channel numbers render on top; reorganize strips accordingly \\
\bottomrule
\end{tabularx}
\vspace{0.5em}

\subsection{Motion Tracking Pitfalls}

\vspace{0.5em}
\rowcolors{2}{rowgray}{white}
\begin{tabularx}{\textwidth}{L{3cm} L{3.5cm} X}
\toprule
\rowcolor{primary}
\tableheader{Pitfall} & \tableheader{Cause} & \tableheader{Solution} \\
\midrule
High solve error & Bad tracking markers or incorrect camera parameters & Remove markers with high individual error; verify focal length is correct \\
Tracking markers slip & Low-contrast features or motion blur in footage & Choose high-contrast corners; increase Pattern Size; use Correlation tracking mode \\
CG objects floating & 3D origin not set correctly after solve & Use Set Origin to place the 3D origin on a tracked ground point \\
Scale mismatch & Scene scale incorrect relative to real world & Use Set Scale on markers with known real-world distance \\
\bottomrule
\end{tabularx}
\vspace{0.5em}

\subsection{Performance Tips}

\begin{itemize}[leftmargin=2em]
  \item \textbf{Compositor:} Use the Viewer node sparingly. Disable Backdrop during complex node tree edits to avoid constant recomputation of the preview.
  \item \textbf{VSE:} Generate \textbf{Proxy} files for smooth playback of high-resolution footage. Set up proxies in Properties $\rightarrow$ Proxy $\rightarrow$ Set Selected Strip Proxies.
  \item \textbf{Motion Tracking:} Track at half resolution first for speed, then refine problem markers at full resolution.
  \item \textbf{Rendering:} For iterative compositing, render to Multi-Layer EXR once, then work from the saved file rather than re-rendering each time you adjust the compositor. This is the single most impactful performance optimization for compositing workflows.
\end{itemize}

\vspace{1em}
\begin{center}
\textcolor{bordergray}{\rule{0.5\textwidth}{0.5pt}}
\end{center}
\vspace{1em}

This chapter has covered the full post-production pipeline available within Blender --- from the node-based Compositor and its render pass system, through the Video Sequence Editor for non-linear editing, motion tracking for VFX integration, and Grease Pencil for 2D/3D hybrid work. The key principle throughout is \textit{non-destructive flexibility}: render once to Multi-Layer EXR, then adjust everything else without re-rendering. Master this workflow, and you will spend your time making creative decisions rather than waiting for renders.

The techniques in this chapter connect directly to the material covered in Chapter~\ref{ch:shading} (where you set up the materials and lighting that produce the render passes), Chapter~\ref{ch:animation} (where you create the motion that the VSE will edit), and Chapter~\ref{ch:furryarts} (where the specific requirements of fur rendering and community content creation inform your compositing choices). The compositor is where all of these threads come together into a finished image.
